\chapter{Введение}
\label{ch:intro}

    В ряде практических задач по проектированию и строительству энергоэффективного жилья необходимо проводить анализ потерь тепла, чтобы оптимизировать затраты на отопление.
    
    При проектировании системы отопления нужно учитывать, как интенсивно воздух будет терять тепло. Основным источником потерь являются окна. Потери тепла через окна зависят от того, каким веществом заполнено пространство между стёклами. 
    
    Можно заполнить это пространство вакуумом, тогда потери тепла будут минимальны. Однако для использования вакуума нужны будут специальные окна, которые смогут выдерживать атмосферное давление на площади порядка единиц квадратных метров. Поэтому этот вариант не подходит.
    
    Второй вариант заключается в заполнении этого пространства воздухом. Этот способ прост в применении. Для того чтобы рассчитать потери тепла и настроить отопление, необходимо произвести расчёт коэффициента теплопроводности воздуха.
    
    Целью работы являлось измерение коэффициента теплопроводности воздуха при атмосферном давлении в зависимости от температуры.
